\documentclass[12pt]{article}

\setlength{\parindent}{0pt}
\setcounter{secnumdepth}{3}
\usepackage[margin=1in]{geometry}

% packages
\usepackage{amsmath, amsfonts, amsthm, amssymb}
\usepackage{graphicx, float}
\usepackage{mathtools}
\usepackage{titlesec}
\usepackage{interval}
\usepackage{hyperref}
\usepackage{siunitx}
\usepackage{titling}
\usepackage{vwcol}
\usepackage{setspace}
\usepackage{empheq}
\usepackage{cancel}
\usepackage{esdiff}
\usepackage{multicol}
\usepackage{mdframed}
\usepackage{esdiff}
\usepackage{tikzsymbols}
\usepackage{multicol}
\usepackage{tikz}
\usepackage{varwidth}
\usepackage{pgfplots}
\usepackage{fancyhdr}
\usepackage{enumitem}
\usepackage{tcolorbox}

\intervalconfig {
	soft open fences
}

% header
\pagestyle{fancy}
\fancyhf{}
\lhead{Ethan Fan}
\rhead{AP Calculus BC}
\cfoot{\thepage}

% section format
\titleformat{\section}
  {\large \bfseries}
  {}
  {0em}
  {}
\titlespacing*{\section}{0pt}{0.5em}{0.3em}

\titleformat{\subsection}[runin]
  {\bfseries}
  {}
  {0em}
  {}
\titlespacing*{\subsection}{0pt}{0pt}{0em}

\titleformat{\subsubsection}[runin]
  {\itshape}
  {(\roman{subsubsection})}
  {0em}
  {}

% commands
\newcommand{\alignedintertext}[1]{%
  \noalign{%
    \vskip\belowdisplayshortskip
    \vtop{\hsize=\linewidth#1\par
    \expandafter}%
    \expandafter\prevdepth\the\prevdepth
  }%
}

\newcommand*{\paren}[1]{\ensuremath\left(#1\right)}
\newcommand*{\limit}[2][x]{\ensuremath{\displaystyle\lim_{#1\to#2}}}
\newcommand*{\Limit}[3][x]{\ensuremath{\displaystyle\lim_{#1\to#2}\left[#3\right]}}
\newcommand*{\deriv}[1][x]{\ensuremath{\dfrac{\mathrm{d}}{\mathrm{d}#1}}}
\newcommand*{\Deriv}[2][x]{\ensuremath{\dfrac{\mathrm{d}}{\mathrm{d}#1}\left[#2\right]}}
\newcommand{\dydx}{\dfrac{dy}{dx}}
\newcommand{\dydt}{\dfrac{dy}{dt}}
\newcommand{\dxdt}{\dfrac{dx}{dt}}
\newcommand*{\iinteg}[2][x]{\ensuremath{\displaystyle\int #2\;\mathrm{d}#1}}
\newcommand*{\dinteg}[4][x]{\ensuremath{\displaystyle\int_{#2}^{#3}#4\;\mathrm{d}#1}}
\newcommand*{\abs}[1]{\ensuremath{\left|#1\right|}}
\newcommand*{\eps}{\varepsilon}
\newcommand*{\real}{\ensuremath\mathbb{R}}
\newcommand*{\integer}{\ensuremath\mathbb{Z}}
\newcommand*{\posint}{\ensuremath\mathbb{Z^+}}
\newcommand*{\cbrt}[1]{\ensuremath{\sqrt[3]{#1}}}
\newcommand*{\lheqzero}{\ensuremath{\underset{\text{L'H}}{\overset{\left[\frac00\right]}{=}}}}
\newcommand*{\lheqinfty}{\ensuremath{\underset{\text{L'H}}{\overset{\left[\frac{\infty}{\infty}\right]}{=}}}}
\newcommand{\tor}{\text{ or }}
\newcommand{\floor}[1]{\lfloor #1 \rfloor}

\newcommand{\vem}{\vspace{0.5em}}
\newcommand{\example}[1]{\section{Example #1}}
\newcommand{\problem}[1]{\section{Problem #1}}
\newcommand{\subproblem}[1]{\subsection{(#1)} \,}

\DeclareMathOperator{\dne}{DNE}
\DeclareMathOperator{\sgn}{sgn}

\newcommand{\definition}[1]{\begin{tcolorbox}[colback=red!5!white,colframe=red!75!black,parbox=false] #1 \end{tcolorbox}}

\newcommand{\theorem}[2]{\begin{tcolorbox}[title={#1},colback=blue!5!white,colframe=blue!75!black,parbox=false] #2 \end{tcolorbox}}

\allowdisplaybreaks
% \postdisplaypenalty=100000

% document
\begin{document}

\vspace*{0cm}

% opening
\begin{center}
    {\Large \bfseries Problem Set 4} \\[0.5cm]
    {\large The Definite Integral} \\[0.35cm]
    {\normalsize September 14, 2026}
\end{center}

% problems

\problem{1}

\subproblem{a}
The drone is climbing from $t\in (0,3)$. The drone is descending from $t\in (3,5)$. It is highest at $t=3$. \vem

\subproblem{b}
\begin{gather*}
    \overline{S}_n=10m\\
    \underline{S}_n=0m\\
\end{gather*}
$\underline{S}_n$ uses right endpoints as the function is monotonically decreasing.\vem

\subproblem{c}
\[
\int_0^5 v(t) \, dt=9-4=5m
\]
The drone's overall displacement in the interval is positive.\vem

\subproblem{d}
\begin{gather*}
    v(t)=6-2t\\
    R_n=\frac{5}{n} \sum_{k=1}^nv\paren{\frac{5k}{n}}=30-\frac{25n(n+1)}{n^2}=5-\frac{25}{n}\\
    \lim_{n \to \infty} R_n=\lim_{n\to \infty} 5-\frac{25}{n}=5
\end{gather*}
The limit is the same as the result in (c).\vem

\subproblem{e}
\begin{gather*}
    \int_0^5 |v(t)| \, dt =9+4=13m
\end{gather*}
It is different from (c) as it considers the sign of the values. The displacement represents the drone's altitude, while distance is the log of meters flown. \vem

\subproblem{f}
\[
M_5=5m
\]
The midpoint sum is exact, as the left side of the midpoint balances with the right side of the midpoint.

\problem{2}

\subproblem{a}
\begin{gather*}
    x(x-2)=0 \implies x=0,2\\
    y=(x-1)^2-1\implies x=1\\
    f(3)=3
\end{gather*}

\subproblem{b}
\begin{gather*}
    R_n=\frac{3}{n} \sum_{k=1}^n f\paren{\frac{3k}{n}}\\
    =\frac{3}{n}(\frac{9}{n^2}\cdot \frac{n(n+1)(2n+1)}{6}-\frac{6}{n} \cdot \frac{n(n+1)}{2})\\
    =\frac{9(n+1)(2n+1)}{2n^2}-\frac{9(n+1)}{n}\\
    =\frac{9(n+1)}{n} (\frac{2n+1}{2n} - \frac{2n}{2n})\\
    =\frac{9(n+1)}{2n^2}\\
    \int_0^3 (x^2-2x) \, dx=\lim_{n \to \infty} R_n =0
\end{gather*}

\subproblem{c}
\begin{gather*}
    R_n=\frac{2}{n} \sum_{k=1}^n f\paren{\frac{2k}{n}}\\
    =\frac{2}{n}(\frac{4}{n^2}\cdot \frac{n(n+1)(2n+1)}{6}-\frac{4}{n} \cdot \frac{n(n+1)}{2})\\
    =\frac{4(2n^2+3n+1)}{3n^2} - \frac{4(n+1)}{n}\\
    =\frac{8}{3}+\frac{4}{n}+\frac{4}{3n^2}-4-\frac{4}{n}\\
    =-\frac{4}{3}+\frac{4}{3n^2}\\
    \int_0^2 (x^2-2x) \, dx = \lim_{n \to \infty} R_n =-\frac{4}{3}\\
    \int_2^3 (x^2-2x) \, dx = \frac{4}{3}\\
    A=\frac{8}{3}
\end{gather*}

\subproblem{d}
\begin{gather*}
    \int_0^3 (x-1)^2 \, dx = \int_0^3 x^2-2x+1 \, dx = 0+3=3\\
    \int_0^3 (x-1)^2 \, dx = \int_{-1}^2 x^2 \, dx =3
\end{gather*}


\problem{5}

\subproblem{a}
\begin{gather*}
    a=-1,b=1\\
    f\paren{\frac{k}{n}}=1-\frac{k^2}{n}\\
    R_n=\frac{2}{n}\cdot \sum_{k=1}^n f\paren{\frac{k}{n}}=\\
    \frac{2}{n}(n-\frac{1}{n^2} \cdot \frac{n(n+1)(2n+1)}{6})\\
    =2-\frac{2n^2+3n+1}{3n^2}\\
    =\frac{4}{3}-\frac{1}{n}-\frac{1}{3n^2}=\boxed{\frac{4}{3}}
\end{gather*}

\subproblem{b}
\begin{gather*}
    \sqrt{1+x} \ge \sqrt{1+x^2}, \quad x\in [0,1]\\
    \int_0^1 \sqrt{1+x} \, dx\ge \int_0^1 \sqrt{1+x^2} \, dx\\
    1 \le \int_0^1 \sqrt{1+x} \, dx \le \sqrt{2}
\end{gather*}

\subproblem{c}
\begin{gather*}
    g(x)=-g(x) \implies \int_{-a}^ag(x) \, dx=0\\
    g(x)=g(-x) \implies \int_{-a}^ag(x) \, dx= 2 \int_{0}^ag(x) \, dx\\
    \int_{-1}^1 x^3\sqrt{1-x^2} \, dx =0\\
    \int_{-1}^1 (x^5+3\sqrt{1-x^2}) \, dx = \frac{3\pi}{2}
\end{gather*}

\subproblem{d}
\begin{gather*}
    \pi \cdot 1\cdot \frac{1}{2}=\frac{\pi}{2}\\
    \pi \cdot \frac{\pi}{2} \cdot \frac{1}{2}=\frac{\pi^2}{4}\\
    \frac{\pi}{2} \le I\le \frac{\pi^2}{4}\\
    W=\frac{\pi^2}{4}-\frac{\pi}{2}=0.90
\end{gather*}

\problem{6}

\subproblem{a}
\begin{gather*}
    R_n=\sum_{k=1}^n \sqrt{x_k} \Delta x\\
    x_k=\frac{k^2}{n^2}\\
    \Delta x=\frac{2k-1}{n^2}\\
    \frac{2n-1}{n^2} \to 0 \implies n\to \infty\\
    R_n=\frac{1}{n^3}\sum_{k=1}^n 2k^2-k\\
    =\frac{1}{n^3}\paren{\frac{n(n+1)(2n+1)}{3}-\frac{n(n+1)}{2}}\\
    =\frac{n+1}{n^2} \paren{\frac{2n+1}{3}-\frac{1}{2}}\\
    \lim_{n\to \infty } R_n=\frac{2}{3}\\
    \boxed{\int_0^1 \sqrt{x} \,  dx =\frac{2}{3}}
\end{gather*}

\subproblem{b}
They are reflections over $y=x$.
\begin{gather*}
    \int_0^1 x^2 \, dx =\frac{1}{3}x^3 \Big|^1_0=\frac{1}{3}\\
    1-\frac{1}{3}=\frac{2}{3}
\end{gather*}

\subproblem{c}
\begin{gather*}
    x_k=b\frac{k^2}{n^2}\\
    \Delta x =b\frac{2k-1}{n^2}\\
    \sqrt{x_k}=\sqrt{b}\frac{k}{n}\\
    R_n=\sum_{k=1}^n \sqrt{x_k} \Delta x=\sum_{k=1}^n \paren{\sqrt{b}\frac{k}{n}}\paren{b\frac{2k-1}{n^2}}\\
    \boxed{\int_0^b \sqrt{x}\, dx= \lim_{n \to \infty} R_n= \frac{2}{3}b^\frac{3}{2}}
\end{gather*}

\subproblem{d}
\begin{gather*}
    \Delta x=\frac{1}{2(n-1)}\\
    R_n=\sum_{k=1}^{n-1} \frac{k}{2n(n-1)}+p_n\cdot \frac{1}{2}=\frac{n+1}{8(n-1)}+p_n\cdot \frac{1}{2},\quad p_n\in [\frac{1}{2},1]\\
    \boxed{\lim_{n\to \infty} R_n=\frac{1}{8}+\frac{p_n}{2}}
\end{gather*}
The $\Delta x$ for every single partition must go to zero. As the $\Delta x$ for the last piece remains $\dfrac{1}{2}$ regardless of the value of $n$, the sums do not converge.




\end{document}

